\documentclass[master, fontsize=12pt]{diploma}

\usepackage[reviewer]{review}

\student{Инютин Максим Андреевич}
\faculty{№ 8 <<Компьютерные науки и прикладная математика>>}
\department{Образовательный центр Института № 8}
\speciality{01.04.02 Прикладная математика и информатика}
\profile{Продуктовая разработка и разработка программных систем}
\group{М8О-206СВ-24}

\theme{Разработка алгоритма лидарной одометрии беспилотного автомобиля, учитывающего перекрытие поля зрения}
% Дата. Оставляем пустое место для дня
\workDate{\uline{\hspace{24pt}} мая \the\year\ года}

\supervisor{Ухов Пётр Александрович}
\supervisorCredentials{кандидат технических наук, доцент, преподаватель кафедры 806}

\firstConsultant{Евсеев Дмитрий Сергеевич}
\firstConsultantCredentials{доктор физико-математических наук, профессор, профессор кафедры 806}

% \secondConsultant{Алешина Елена Константиновна}
% \secondConsultantCredentials{доктор физико-математических наук, профессор, профессор кафедры 806}

% \reviewer{Свиридов Михаил Артемьевич}
% \reviewerCredentials{доктор физико-математических наук, профессор}

\headOfDepartmentTitle{И. о. директора Образовательного центра}
\headOfDepartment{Крылов Сергей Сергеевич}



\begin{document}
\fillReviewHeader

\textbf{Отмеченные достоинства:}

Выпускная квалификационная работа Инютина Максима Андреевича посвящена актуальной задаче лидарной одометрии беспилотного автомобиля при перекрытии поля зрения. Тема практически значима для автономного вождения: крупные динамические объекты ухудшают наблюдаемость сцены и могут приводить к деградации классических методов регистрации облаков точек.

К достоинствам работы следует отнести чёткую постановку задачи: требуется не длительное движение только по одометрии, а устойчивое прохождение короткого критического участка до восстановления большей наблюдаемости сцены. Такая постановка хорошо согласована с ограничениями GNSS и особенностями лидарных данных.

Автор предложил и реализовал устойчивый алгоритм, включающий предобработку и двойную вокселизацию облаков точек, ICP с оценкой преобразования между парами точек с учётом весов, модель постоянной скорости и фильтр Калмана. Функция Гемана---Макклюра снижает влияние выбросов, а фильтр Калмана стабилизирует начальное приближение следующего шага ICP. Такое сочетание методов инженерно обосновано.

Экспериментальная часть выполнена на данных Boreas-RT: 10 тестовых проездов разбиты на 725 десятисекундных сцен с моделируемым обгоном грузовика. Результаты приведены на характерных примерах и в агрегированной статистике по основным метрикам ошибок; значимость различий проверена парным критерием Стьюдента. Это повышает убедительность выводов и воспроизводимость сравнения.

Результаты показывают, что подход заметнее всего улучшает оценку продольной скорости, особенно чувствительной к ухудшению наблюдаемости. При этом фильтр Калмана не подавляет естественную динамику движения на длинных последовательностях, а снижает риск последовательной деградации регистрации. Важным достоинством является честное обсуждение ограничений метода, включая геометрически вырожденные сцены в туннелях и на трассовых участках.

\textbf{Отмеченные недостатки:} Параметры ковариаций процесса и измерения в фильтре Калмана выбраны по инженерным соображениям и не оптимизируются по экспериментальным данным. Работа также могла бы выиграть от более подробного сравнения с картографическими методами лидарной локализации или подходами с дополнительными источниками информации. Эти замечания носят рекомендательный характер и не снижают общей положительной оценки.

\textbf{Заключение}: Выпускная квалификационная работа выполнена на высоком уровне, содержит актуальную постановку задачи, самостоятельную программную реализацию, воспроизводимый экспериментальный конвейер и содержательный анализ результатов. Работа соответствует требованиям к магистерской диссертации. Инютин Максим Андреевич рекомендуется к присвоению квалификации <<инженер-исследователь>> по направлению 01.04.02 <<Прикладная математика и информатика>>. Рекомендуемая оценка --- <<отлично>>.

\makeReviewSignature
\end{document}
